\section{Structured case-search record}\label{si:cases}

\subsection{Eligibility criteria and adjudication rule}

The inventory contains 32 candidate systems. It is a curated query-driven search, not a systematic review. A candidate was eligible only if all four criteria were supported by evidence:
\begin{enumerate}
\item an operative, repeatedly responsive governance rule (not a one-time cap, ban, fixed path, coordination deadlock, or nominal rule with no effective response);
\item an independently dateable response or implementation lag at the scale relevant to the ecological outcome;
\item no competing ecological or structural driver supported for that same outcome, or comparative evidence that institutional feedback predicts the outcome better than prespecified competing mechanisms; and
\item individual-resource or station-level observations adequate to order the response signal, decision, implementation, and ecological effect.
\end{enumerate}
The third condition is not satisfied merely because an institutional mechanism is plausible. If ecology, recruitment, climate, predation, cohort resonance, or assessment structure is supported for the same outcome and no discriminating comparison exists, the case is coded as confounded or unresolved. ``Yes'' in the table means that the available record supports the item; ``no'' means the criterion fails; a dot denotes no determination or inapplicability; ``partial'' denotes incomplete evidence. The binding-fail column is the recorded failure, not a composite score.

Twenty cases are coded as having a responsive rule, with one additional case coded partial; 21 rows record criterion (iii) as unmet. Two produced-capital examples are out of the regenerating-resource scope. No row satisfies all four criteria simultaneously. The closest candidates—Bangkok groundwater and La Mancha Oriental groundwater—remain confounded by subsidence/rainfall and crop/water substitution, respectively. These decisions should be read as limits of this inventory, not as proof that no qualifying system exists.

\begin{landscape}
\scriptsize
\setlength{\tabcolsep}{3pt}
\begin{longtable}{@{}p{0.55cm}p{3.3cm}p{1.6cm}p{1.0cm}p{1.6cm}p{1.05cm}p{1.0cm}p{1.0cm}p{10.8cm}@{}}
\caption{Curated 32-system inventory and four-criterion dispositions. A dot denotes no assessment or not applicable; the binding-fail column records the criterion or scope exclusion that prevents full eligibility.}\label{tab:case-inventory}\\
\toprule ID & System & Domain & C1 & C2 lag & C3 & C4 & Binding fail & Evidence note \\
\midrule\endfirsthead
\toprule ID & System & Domain & C1 & C2 lag & C3 & C4 & Binding fail & Evidence note \\
\midrule\endhead
\bottomrule\endfoot
1 & Baltic sprat (SPRAT22-32 EU MAP) & fisheries & yes & yes (\textasciitilde{}1-1.5 yr) & no & yes & iii & Annual TAC responsive; variance is low-frequency trend/regime (no robust cycle); tau-window screen row; RAM SSB series \\
2 & European anchovy Adriatic (GFCM MAP) & fisheries & no & .  & .  & .  & i & 2022-2026 fixed predetermined reduction path: no deficit feedback (GFCM/44/2021/20; GFCM/48/2025/2) \\
3 & European anchovy Bay of Biscay & fisheries & yes & yes (\textasciitilde{}1-1.5 yr) & no & yes & iii & Annual TAC; recruitment/environment-driven variability; closed 2005-2009 \\
4 & Peruvian anchoveta & fisheries & yes & yes (sub-annual) & no & yes & iii & ENSO driver: 3.7-yr peak Granger-caused by ENSO (section 3.7 battery) \\
5 & Pacific sardine (US\allowbreak/Mexico) & fisheries & yes & yes (\textasciitilde{}0.5-1 yr) & no & yes & iii & Annual HCR; environmentally forced regime variability \\
6 & Atlantic menhaden (ASMFC) & fisheries & yes & yes (\textasciitilde{}1-1.5 yr) & no & yes & iii & Annual quota; environmentally forced recruitment variability \\
7 & Gulf menhaden (GSMFC) & fisheries & yes & yes (\textasciitilde{}1-1.5 yr) & no & yes & iii & Annual quota (5-state commission); environmentally forced variability \\
8 & North Sea herring (ICES) & fisheries & yes & yes (\textasciitilde{}1-1.5 yr) & no & yes & iii & Annual TAC; recruitment/climate confounds \\
9 & Atlanto-Scandian herring & fisheries & no & .  & .  & .  & i & Coastal-state sharing deadlock since \textasciitilde{}2014; unilateral quotas above advice: coordination failure not a rule \\
10 & NE Atlantic mackerel \allowbreak/ blue whiting & fisheries & no & .  & .  & .  & i & Coastal-state deadlock (MSC suspended 2019-2020): no operative responsive rule \\
11 & Iberian sardine (PT\allowbreak/ES) & fisheries & yes & yes (\textasciitilde{}1-1.5 yr) & no & yes & iii & Annual TAC + rebuilding measures; environmentally forced variability \\
12 & South Africa anchovy\allowbreak/sardine (OMP) & fisheries & yes & yes (\textasciitilde{}1-1.5 yr TAC) & no & yes & iii & Annual TAC within OMP (rule review 4-5 yr is not effort-response lag); environmentally forced \\
13 & Namibian sardine\allowbreak/anchovy & fisheries & yes & yes (\textasciitilde{}1-1.5 yr) & no & yes & iii & Annual TAC; environmentally forced variability \\
14 & Chilean anchovy \allowbreak/ jack mackerel & fisheries & yes & yes (\textasciitilde{}1 yr) & no & yes & iii & Annual TAC + seasonal management; environmentally forced variability \\
15 & Japanese sardine & fisheries & yes & yes (\textasciitilde{}1 yr) & no & yes & iii & Annual TAC; regime-shift-driven variability \\
16 & West African sardinella & fisheries & no & no & .  & .  & i,ii & Annual TACs exist but enforcement/implementation weak and irregular (CECAF): no dateable responsive lag \\
17 & NZ QMS pelagics & fisheries & no & .  & .  & .  & i & TACs held between sparse 3-10 yr assessments (57\% never changed): discontinuous not delayed-continuous (Mace et al. 2013/2014) \\
18 & Barents \allowbreak/ Icelandic capelin & fisheries & yes & yes (<1-1 yr) & no & yes & iii & Annual quotas + in-season surveys; dynamics cod-driven/environmentally forced \\
19 & North Sea sandeel & fisheries & yes & yes (\textasciitilde{}1-1.5 yr) & no & yes & iii & Annual TAC seasonal; environmentally forced + seabird-driven closures \\
20 & India\allowbreak/Philippines\allowbreak/Indonesia small pelagics & fisheries & no & .  & .  & .  & i & No TAC; seasonal closures/mesh regulations only: fails responsive-feedback criterion \\
21 & Icelandic cod (1995 HCR) & fisheries & yes & yes (0.2-0.3 yr impl. lag) & no & yes & iii & 25\% HCR responsive; cohort resonance alternative (section 3.7/S4); post-rule SSB CV 0.387 (audit-confirmed) \\
22 & Icelandic haddock (2013 HCR) & fisheries & yes & yes (\textasciitilde{}1 yr TAC cycle) & no & yes & iii & 35\%-of-B45cm+ HCR adopted Apr 2013 (Iceland govt); recruitment-forced variability (R CV 0.89) \\
23 & Sheridan-6 groundwater & groundwater & no & .  & .  & yes & i & 55 acre-inch block allocation over five-year first period: not continuously adjusted feedback (S4) \\
24 & Bangkok groundwater & groundwater & yes & yes (post-1999 response) & no & yes & iii & Pumping declines after 1999 but subsidence/recovery confounded by heterogeneous consolidation + rainfall (S4) \\
25 & La Mancha Oriental groundwater & groundwater & . & . & no & . & iii & 2019-2023 extraction relapse; attribution complicated by crops/ET/surface-water substitution (S4) \\
26 & Murray-Darling surface water & surface water & yes & yes (annual announcements) & no & yes & iii & Allocations responsive to storage/inflows (SA DEW; BoM NWA 2024) but availability is rainfall-driven: climate confound \\
27 & US BLM rangeland grazing & rangeland & partial & no & .  & . & ii & Responsive adjustments only via recent flexibility pilots (IM 2018-109; outcome-based grazing); historically rigid permits with NEPA-gated irregular lags: no dateable response lag \\
28 & Scandinavian moose harvest & wildlife & yes & yes (annual/3-5 yr plans) & no & yes & iii & Responsive quotas (SE annual county; NO 3-5 yr plans) but recruitment climate-driven (hot dry springs/snow) + wolf predation confounds \\
29 & Norwegian salmon aquaculture (TLS) & aquaculture & yes & yes (dated evaluations 2017+) & no & yes & iii & Biomass responsive to lice-induced wild-smolt mortality (TLS 13 areas) but lice signal biologically/environmentally forced + farm economics joint: institutional channel not isolated \\
30 & Spruce budworm forestry & forestry & no & .  & .  & . & i & 30-40 yr cycle is endogenous predator-prey outbreak dynamics (Ludwig Jones Holling 1978): no institutional feedback at all \\
31 & Hog\allowbreak/livestock cycles & produced capital & . & . & . & . & scope & Unconfounded delay oscillators but no autonomously regenerating stock: outside the resource-model scope (section 3.7/S4) \\
32 & Electricity-capacity cycles & produced capital & . & . & . & . & scope & Unconfounded delay oscillators but no autonomously regenerating stock: outside the resource-model scope (section 3.7/S4) \\
\end{longtable}
\end{landscape}


\subsection{Search provenance and query-log reconciliation}

The retained \texttt{\detokenize{u5_query_log.md}} is dated 2026-09-11. It records (i) repository tree and code searches; (ii) targeted extraction of relevant release files; (iii) RAM Legacy v4.66 re-execution and the cod/haddock data audit; (iv) ICES stock-data pulls for northern cod and haddock; and (v) primary-source web searches for water allocation, grazing, moose harvest, salmon-lice governance, and Icelandic harvest-control rules. The table's first 20 rows were re-adjudicated under the four criteria rather than mechanically inherited from an earlier delay-window screen. Row-level notes and search queries are preserved verbatim in the source bundle.

The query log's search R4 found a memo that called the \(\Tr\)-ranked test ``never executed.'' That was a historical snapshot of the analysis status, not the final repository state: the executable test script and outputs were subsequently recovered and re-run. The discrepancy is documented in S1 and the source reconciliation note; it does not change the case-screening result.

